\begin{savequote}
\includegraphics[width=5cm]{./images/teaser/follow-me-kontur.jpg}
\end{savequote}

\chapter
[\CO{[PAM]} Patientenmanagement]
{\CC{[PAM]} Patientenmanagement}

\minitoc

\includegraphics[angle=270,width=\linewidth]{./images/leitsystem/Leitschema-PAM_gesamt.pdf}

\Layout{2}{Inhalt}{%
Dieses Kapitel behandelt den grundsätzlichen Ablauf eines
Einsatzes und der Versorgung von Patienten. Grundbausteine sind dabei die
Kenntnis der richtigen Reihenfolge der Maßnahmen und die Fähigkeit die
Übersicht zu behalten.  Das Erarbeiten von (Verdachts-)Diagnosen mittels
Patientengespräch und Untersuchungen sowie das Beherrschen der theoretischen
medizinischen Grundlagen sind weitere Voraussetzungen für die erfolgreiche
Patientenbetreuung und effektives, professionelles Arbeiten.
}{}{}{}{}


\Layout{2}{Querverweise}{%
\begin{compactdesc}
  \item[Beschreibung der Medizinprodukte:] Kapitel [MP1, MP2] 
  (\ref{chp:MP1}, \ref{chp:MP2}).
  \item[Untersuchungen:] Kapitel [BAU] (\ref{chp:BAU}).
  \item[Sanitätshilfemaßnahmen:]  Kapitel [SAN] (\ref{chp:SAN}).
  \item[Anamnese und Patientengespräch:]  Kapitel [AUP] (\ref{chp:AUP}).
\end{compactdesc}
}{}{}{}{}



\section{"`Hilfe, ein Patient!"' - Was mach{\textquotesingle} ich jetzt?}

Gerade zu Beginn der Tätigkeit als medizinische Fachkraft
fühlt man sich beim Eintreffen am Notfallort durch die große Menge an Sinneseindrücken und
Informationen überlastet. Um Sie sicher durch diese Informationsflut zu
leiten, wollen wir Ihnen in diesem Abschnitt ein paar Anhaltspunkte
präsentieren, welche den Einsatzablauf strukturieren und Ihnen -- und somit auch
dem Patienten -- Sicherheit geben.

\Layout{0}{Vorgehen am Notfallort:}{%
\begin{itemize}
  \item Bevor es los geht: Szeneüberblick und Selbstschutzmaßnahmen
  \item 1. Schritt: \Es{Einschätzung des Patienten}
  \item 2. Schritt: \Es{Erstmaßnahmen}
  \item 3. Schritt: \Es{Durchatmen, nachdenken, planen}
\end{itemize}
}{
\begin{itemize}
  \item Szeneüberblick und Selbstschutzmaßnahmen
  \item Einschätzung des Patienten
  \item Erstmaßnahmen
  \item Durchatmen, nachdenken, planen
\end{itemize}
}{}{}{}


Während die ersten zwei Schritte ein gezieltes Vorgehen erfordern, bietet der
dritte Schritt Spielraum für die individuelle Untersuchung und Befragung des
Patienten entsprechend den Symptomen. Im Folgenden werden die Schritte
im Einzelnen vorgestellt.

% \clearpage
\section{Vorgehen in der Praxis}

\subsection{Primäreinsatz}

%\paragraph{Szeneüberblick}
% \paragraph{Die ersten  30 Sekunden }

% \Fazit{Ersteinschätzung ${\rightarrow}$~Erstmaßnahmen
% ${\rightarrow}$~Plan zurechtlegen ${\rightarrow}$~Weitere Untersuchung,
% Anamnese und Versorgung}





\begin{gWideParEnv}
\aBox{Ablauf}{Patientenmanagement}{Beige}
{\sffamily
\begin{tabularx}{\linewidth}{cX}
\TabHeading{Einschätzung}
&
\noindent\fcolorbox{green!10}{green!10}{
% \begin{minipage}{\linewidth}
\begin{tabular}{lll}
\toprule\addlinespace
\TabSub{}
&
&
\TabSub{Szeneüberblick}
\\\addlinespace\midrule\addlinespace
\TabSub{\Z{GI}}
&
\TabSub{\Z{General Impression}}
&
 \TabSub{Eindruck}
\\\addlinespace
\TabSub{\Z{}}
&
&
 \TabSub{Bewusstsein}
\\\addlinespace
\TabSub{\Z{}}
&
\TabSub{\Z{}}
&
\TabSub{Hauptbeschwerde}
\\\addlinespace\midrule\addlinespace
\TabSub{\Z{A}}
&
\TabSub{\Z{Airway}}
&
 \TabSub{Atemweg}
\\\addlinespace
\TabSub{\Z{B}}
&
\TabSub{\Z{Breathing}}
&
\TabSub{Atmung}
\\\addlinespace
\TabSub{\Z{C}}
&
\TabSub{\Z{Circulation}}
&
 \TabSub{Kreislauf}
\TabSub{\Z{}}
\\\addlinespace
\TabSub{\Z{D}}
&
\TabSub{\Z{Disability}}
&
\TabSub{Schnelle Untersuchung}
\\\addlinespace\bottomrule
\end{tabular}
% \end{minipage}
}
\\\addlinespace
${\downarrow}$
&
\\\addlinespace
\noindent\TabHeading{Erstmaßnahmen}
&
Je nach Ergebnis der Einschätzung. 

Z.B. bei vitaler Bedrohung: \TxMassVit.
\\\addlinespace
${\downarrow}$
&
\\\addlinespace
\noindent\TabHeading{Plan zurechtlegen} 
&
\\\addlinespace
${\downarrow}$
&
\\\addlinespace
\noindent\TabHeading{Weitere Untersuchung,
Anamnese und Versorgung}
&
Übersicht Tafel \prettyref{tab:checkliste-patientenmanagement}
\\\addlinespace
\end{tabularx}




% \sffamily 
% \begin{center}
% Ersteinschätzung 
% 
% ${\downarrow}$
% 
% Erstmaßnahmen
% 
% ${\downarrow}$
% 
% Plan zurechtlegen 
% 
% ${\downarrow}$
% 
% Weitere Untersuchung,
% Anamnese und Versorgung
% \end{center}
}
\end{gWideParEnv}



\subsubsection{Erster Schritt: Einschätzung}
\label{topic:ersteinschaetzung}

\Layout{60}{Szeneüberblick}{%
Schon \E{vor} Eintreffen am Einsatzort beginnt der Szeneüberblick. Das
Einsatzpersonal muss sich die äußeren Umstände des Einsatzes bewusst
machen: \Es{Wo} befindet sich der Einsatzort, in einer Wohnung in einem
mehrstöckigen Wohnhaus, in einem Einfamilienhaus oder auf einer Verkehrsfläche?
Wo befindet sich der \Es{Eingang}, wo gibt es \Es{Parkmöglichkeiten}? Gibt es
einen \Es{Aufzug} und wie groß ist er? Können \Es{nachfolgende Kräfte} den
Einsatzort leicht finden? 

Beim Eintreffen am Einsatzort ist besonders auf \Es{Gefahren} zu achten,
z.\,B. \E{"`Vorsicht vor dem Hund"'}-Schilder, \E{Futternäpfe}, Hinweise auf
\E{Gasunfälle}, unsichere Gebäude(-teile), mögliche \E{Fluchtwege} usw. Wichtig
ist auch die Einschätzung ob es sich um eine \Es{Großschadenslage} handelt. 

Bei Unfällen gibt die Inspektion des Einsatzortes oft auch Aufschluss über den
\Es{Unfallmechanismus} und -hergang. 
}{
\begin{itemize}
  \item Zufahrt, Gebäude, Gefahren; Großschaden?;
  Umgebung, Ort, Zeit; Umstände, Hinweise auf Unfallmechanismus
\end{itemize}
}
{./images/motal-aass/00800/helm1.jpg}
{}
{}

\Layout{60}{GI -- Eindruck}{%
\Z{(\Lang{engl.} General Impression.)} Hier wird v.\,a. das 
\Es{Aussehen}, der Gesichtsausdruck, die \Es{Körperhaltung} und die \Es{Sprache}
des Patienten beurteilt.

Auffallen kann eine \E{blasse  Haut}, \E{Zyanose}, \E{Schweiß},
\E{Schonhaltung},
\E{Atemnot} (Gestik, Aufstützen der Arme um die \E{Atemhilfsmuskulatur}
einzusetzen (\prettyref{topic:atemhilfsmuskulatur}), Nasenflügeln, Einziehungen an den
Rippen, \ldots{}), eine verwaschene oder wirre \E{Sprache},
\E{Des\-orientiertheit}, \E{Krämpfe}, u.\,v\,a.\,m. 
}{%
\begin{itemize}
  \item Aussehen, Gesichtsausdruck, Haltung, Sprache
  \item Blässe, Zyanose, Schweiß, Schonhaltung, Atemnot, verwaschene
Sprache, Desorientiertheit, Krämpfe 
\end{itemize}
}
{./images/teaser/imgp0722_00500.jpg}
{}
{}

\Layout{0}{GI -- Bewusstsein}
{%
Beurteilt wird der \Es{quantitative
\E{Bewusstseinsgrad}} (\Speech{"`\Es{Wieviel} Bewusstsein ist vorhanden?"'},
vgl. \prettyref{topic:bewusstseinsstoerung-grade}). Es wird dazu die Reaktion auf verbale
Ansprache, bzw. wenn notwendig auf einen Weckversuch oder Schmerzreiz bewertet.
Der Patient kann \E{klar}, \E{somnolent}, \E{soporös} oder \E{komatös}/\E{bewusstlos} sein. (Quantitavive
Bewusstseinsgrade, vgl. \prettyref{topic:bewusstsein}) 
}{%
\begin{itemize}
  \item Reaktion auf verbale Ansprache, Weckversuch, Schmerzreiz
  \item klar, somnolent, soporös, komatös/bewusstlos
\end{itemize}
}{./icons/dead.png}{}{}

\Layout{0}{GI -- Hauptbeschwerde}{%
Die \Es{Hauptbeschwerde} des Patienten, bzw. der \Es{Grund der Berufung}, sowie 
Schmerzen und andere für den Patienten dringliche Symptome müssen herausgefunden
werden. I.\,d.\,R. wird hier auf den "`S"'-Teil des \Abk{S-KAMEL} vorgegriffen.
Es bietet sich die Frage an: 

\begin{quote}
\Speech{"`Was können wir für sie tun?"'}
\end{quote}

Eventuell ist auch die Umgebung (Angehörige, Dritte) einzubeziehen
(Fremdanamnese).
}{%
\begin{itemize}
  \item Eigene Wahrnehmung, Patientengespräch, Angehörige.
  \item Grund der Berufung, Schmerzen und andere für den Patienten dringliche
Symptome.
\end{itemize}

}{./icons/crying.png}{}{}

\Layout{60}{A -- Atemweg}{%
\Z{(\Lang{engl.} Airway.)}
Erhoben bzw. beobachtet werden \Es{Atemgeräusche} und die \Es{Gestik} des
Patienten.

Die Atemwege können \E{frei} oder \E{verlegt} sein.
\vspace{1.4cm}
}{%
\begin{itemize}
  \item Atemgeräusch, Gestik.
  \item frei oder verlegt.
\end{itemize}
}
{./images/hirtler-aass/Fremdkoerper.pdf}
{}
{Hirtler}

\Layout{60}{B -- Atmung}{%
\Z{(\Lang{engl.} Breathing.)}
Untersucht wird die \Es{Frequenz}, das \Es{Atemmuster}, die \Es{Anstrengung} des
Patienten beim Atmen, andere als die bei A beschriebenen Atemgeräusche und das Vorhandensein
von \E{Heimsauerstoff}.

\Z{\E{Speziell ausgebildetes} Personal kann mittels eines Stethoskops die oberen
und unteren Atemwege abhören (auskultieren) und damit Störungen des Atemweges
und der Atmung besser einschätzen.}

}{%
\begin{itemize}
  \item Frequenz, Atemmuster, Anstrengung beim Atmen, Atemgeräusche,
  Heimsauerstoff.
\end{itemize}


}
{./images/hirtler-aass/Atmung.pdf}
{}
{Hirtler}

\Layout{60}{C -- Kreislauf}{%
\Z{(\Lang{engl.} Circulation.)}
Die erste Einschätzung des Kreislaufes dauert nur wenige Sekunden.
Beurteilt wird die \Es{Hautfarbe} und \Es{-temperatur}, ob und wie gut
der \Es{Radialispuls} zu tasten ist, sowie die \Es{Rekap-Zeit} an den
Fingerkuppen. 
}{%
\begin{itemize}
  \item Hautfarbe und -temperatur, Radialispuls, Rekap-Zeit.
\end{itemize}
}
{./images/hirtler-aass/Kreislauf-Schema.pdf}
{}
{Hirtler}

\Layout{2}{D -- Schnelle Untersuchung}{%
\Z{(\Lang{engl.} Disability.)} Hier unterscheidet man im
Untersuchungsablauf zwischen Traumapatienten und Nicht-Traumapatienten: 

\begin{labeling}[]{Nicht-Traumapatienten:}
\item[Bei Traumapatienten:] \begin{compactdesc}
    \item[1. Helfer:] Vitalwerte, 
    \item[2. Helfer:] Schneller Traumacheck.
% \item Suche nach: Schwere Verletzungen, Becken-, Oberschenkel-,
% Serienrippenfraktur, massive Blutung, hartes Abdomen.
\end{compactdesc}
\item[Nicht-Traumapatienten:] \begin{itemize}
    \item Vitalwerte.
\end{itemize}
\end{labeling}
}{

}{}{}{}



\begin{gWideParEnv}
\MassnahmenNG{Standardmaßnahmen}{Einschätzungsblock}{\label{m:ersteinschaetzung}}{
\noindent\begin{minipage}[t]{\linewidth -
{2\fboxrule} - {2\fboxsep}}
    	\noindent\fcolorbox{green!10}{green!10}{\begin{minipage}[t]{\linewidth}
    	{\vspace{-0.8em} 
% \begin{tabulary}{\linewidth}[H]{cp{8em}p{8.5cm}}
\begin{tabulary}{\linewidth}[H]{llL}
\toprule\addlinespace
\TabSub{}
&
\TabSub{Szeneüberblick}
&
Sicherheit, Selbstschutz; Großschaden?; Umgebung, Ort, Zeit; Umstände,
Hinweise auf Unfallmechanismus
\\\addlinespace\midrule\addlinespace
\TabSub{\Z{GI}}
&
 \TabSub{Eindruck}
 &
Aussehen, Gesichtsausdruck, Haltung, Sprache

Blässe, Zyanose, Schweiß, Schonhaltung, Atemnot, verwaschene
Sprache, Des\-orientiertheit, Krämpfe 
\\\addlinespace
\TabSub{\Z{}}
&
 \TabSub{Bewusstsein}
 &
Reaktion auf verbale Ansprache, Weckversuch, Schmerzreiz

klar, somnolent, soporös, komatös/bewusstlos
\\\addlinespace
\TabSub{\Z{}}
&
\TabSub{Hauptbeschwerde}
&
Grund der Berufung, Schmerzen und andere für den Patienten dringliche
Symptome
\\\addlinespace\midrule\addlinespace
\TabSub{\Z{A}}
&
 \TabSub{Atemweg}
 &
Atemgeräusch, Gestik

frei oder verlegt
\\\addlinespace
\TabSub{\Z{B}}
&
\TabSub{Atmung}
&
Frequenz, Atemmuster, Anstrengung beim Atmen, Atemgeräusche,
Heimsauerstoff
\\\addlinespace
\TabSub{\Z{C}}
&
 \TabSub{Kreislauf}
\TabSub{\Z{}}
&
Bewusstsein, Hautfarbe, Radialispuls, Rekap-Zeit
\\\addlinespace
\TabSub{\Z{D}}
&
\TabSub{Schnelle Untersuchung}
&
\begin{minipage}[t]{\linewidth}
\protect\Parallel{
\Es{Bei Traumapatienten:} 
\begin{itemize}
    \item 1. Helfer: Vitalwerte 
    \item 2. Helfer: Schneller Traumacheck
% \item Suche nach: Schwere Verletzungen, Becken-, Oberschenkel-,
% Serienrippenfraktur, massive Blutung, hartes Abdomen.
\end{itemize}
}{
\Es{Bei Nicht-Traumapatienten:} 

\begin{itemize}
    \item Vitalwerte.
\end{itemize}
}
\end{minipage}
\\\addlinespace\bottomrule
\end{tabulary}

    	}
      		\end{minipage}}
		\end{minipage}

}
\end{gWideParEnv}



\vspace{1em}

\noindent Nun stellen sich folgende Fragen:

\begin{enumerate}
  \item \Speech{Was ist das Problem des Patienten?}
  \item \Speech{Ist der Patient akut vital bedroht?} (Alarmzeichen, s.u.)
  \item \Speech{\ldots oder sind noch weitere Untersuchungen zur Lagebeurteilung
  nötig? (z.B. Vitalwerte)}
  \item \Speech{Welche weiteren Untersuchungen, Maßnahmen und Fragen stehen an
  erster Stelle?}
\end{enumerate}

    	

\Layout{60}{Vitale Bedrohung -- Alarmzeichen
{\color{DefaultRed}\bfseries{\VarFlag ~}}} {%
Das rasche Feststellen einer vitalen Bedrohung ist besonders wichtig. Anzeichen
einer vitalen Bedrohung können u.\,a. sein: Störungen von Be\-wusstsein, Atemweg,
Atmung oder Kreislauf sowie typische Symptome wie sich nicht bessernde Atemnot,
brodelndes Atem\-ge\-räusch, Thoraxschmerz, Schocksymptome und entgleiste
Vital\-werte. Manche Diagnosen sind automatisch mit einer vitalen Bedrohung
gleichzusetzen, auch wenn eventuell Symptome fehlen oder nur schwach ausgeprägt
sind.

\Es{Sobald eine vitale Bedrohung erkannt wird, müssen die \TxMassVit{}
 (\prettyref{m:standardmassnahmen-vital}) \E{frühzeitig}
ergriffen werden (evt. sogar noch vor Abschluss des
Einschätzungsblockes)!}}{\vspace{1em}

\fcolorbox{DefaultRed}{White}{
\begin{minipage}[t]{\linewidth - 4\Korrektur}
{
\begin{itemize}
    \item Massive Störungen von Vitalfunktionen
    \item Alarm-Symptome
% \begin{itemize}
%     \item Atemnot, die sich nicht bessert 
%     \item Brodelndes Atemgeräusch 
%     \item Thoraxschmerz  
%     \item Schocksymptome 
%     \item entgleiste Vitalwerte
%     \item schwere Verletzungen
%     \item {\dots} 
% \end{itemize}
\item Alarm-Diagnosen
\end{itemize}

Übersicht Tafel \prettyref{tab:alarmzeichen}

% \begin{itemize}
%     \item Bewusstsein
%     \item Atemweg / Atmung 
%     \item Kreislauf 
% \end{itemize}




}
\end{minipage}
}
}{./images/teaser/minen_hinweis2.jpg}
{Vorsicht!}
{GABS.}


% \marginpar{~\newline}
% \Parallel{
% \fcolorbox{ubuntured}{White}{
%       		\begin{minipage}[t]{\linewidth - 2\Korrektur}
%               \Ca{Alarmzeichen}
%               
%       			{\sffamily\Es{Massive Störungen von} 
% 
% Bewusstsein
% 
% Atemweg / Atmung 
% 
% Kreislauf 
% 
% \Es{Typische Symptome} 
% 
% Atemnot, die sich nicht bessert 
% 
% Brodelndes Atemgeräusch 
% 
% Thoraxschmerz  
% 
% Schocksymptome 
% 
% entgleiste Vitalwerte
% 
% schwere Verletzungen
% 
% {\dots} }
% \end{minipage}
% }}
% {
% \paragraph{Alarmzeichen}
% 
% Anzeichen einer vitalen Bedrohung können u.\,a. sein: Störungen
%         von Be\-wusstsein, Atemweg, Atmung oder Kreislauf sowie typische
%         Symptome wie sich nicht bessernde Atemnot, brodelndes
%         Atem\-ge\-räusch, Thoraxschmerz, Schocksymptome und entgleiste
%         Vital\-werte.
%         
% \begin{center}
% \includegraphics[width=4cm]{./images/teaser/minen_hinweis2.jpg}
% \end{center}
% }

\Layout{57}{Alarmzeichen {\color{DefaultRed}\bfseries{\VarFlag ~}} --
Übersicht}{}{} {
\begin{tabularx}{\linewidth}{XXX}
\TabHeading{Massive Störung}
&
\TabHeading{Alarm-Symptome}
&
\TabHeading{Alarm-Diagnosen}
\\\midrule
Bewusstsein

Atemweg / Atmung 

Kreislauf 
&
Atemnot, die sich nicht bessert 

Brodelndes Atemgeräusch 

Thoraxschmerz  

Schocksymptome 

entgleiste Vitalwerte

schwere Verletzungen

{\dots} 
&
Herzinfarkt

Kardiales Lungenödem

Status Epilepticus

Beckenfraktur

\ldots
\\\bottomrule
\end{tabularx}
}{{\color{DefaultRed}\bfseries{\VarFlag}} \Ca{Alarmzeichen:} Wann ist eine vitale Bedrohung
wahrscheinlich? Eine Übersicht.\label{tab:alarmzeichen}}{}



\marginpar{\scriptsize\Hand{Bei der Prüfung "`blind"' auf
Nummer sicher gehen und alle gleich für "`vital bedroht"'
erklären gilt nicht!}

Viele Patienten bei der Prüfung sind nicht vital bedroht und brauchen
keinen Notarzt. 

Erklärt man einen Patienten für vital bedroht muss man das
begründen können!
}
\noindent Während des weiteren Verlaufs (Anamnesegespräch, etc.) kann man
auf weitere Befunde stoßen welche auf eine vitale Bedrohung hinweisen
(zB. massiv erhöhter Blutdruck bei vorgeschädigtem Herzen,
plötzliche oder schleichende Zustandsverschlechterung, etc.). 

Die Ersteinschätzung muss ständig überprüft werden. Neue Befunde
und der Zustand des Patienten können die Situation ändern.

\subsubsection{Zweiter Schritt: Erstmaßnahmen}
\label{topic:patientenmanagement-erstmassnahmen}

Welche Maßnahmen sind für den Patienten am wichtigsten und
müssen zuerst erfolgen? 








\Layout{2}{Vital bedrohte Patienten}{%
Bei vital bedrohten Patienten gibt es 
\Es{5 {\TxMassVit}} (\prettyref{m:standardmassnahmen-vital}, Seite 
\pageref{m:standardmassnahmen-vital}),
die \E{sobald als möglich} (evt. sogar
noch vor Abschluss des Ersteinschätzungsblockes) zu erfolgen haben (
\begin{inparaenum}
\item Situationsgerechte Lagerung, 
\item Sauerstoffgabe (sofern keine Kontraindikation), 
\item Notarztnachforderung (mit Ausnahmen), 
\item Reanimationsbereitschaft, 
\item Engmaschiges, bestmögliches Monitoring.
\end{inparaenum}
).
}{

}{}{}{}



\Layout{2}{Nicht vital bedrohte Patienten}{
Bei
nicht vital bedrohten Patienten muss individuell überlegt werden welche
Maßnahmen zuerst zu treffen sind. 
}{

}{}{}{}


\Layout{2}{Immer durchzuführende Standardmaßnahmen}{%
\prettyref{m:standardmassnahmen-immer} auf Seite
\pageref{m:standardmassnahmen-immer} listet Maßnahmen auf welche bei jeder
Patientenbetreuung -- sofern angemessen -- durchgeführt werden müssen.  }{

}{}{}{}


\subsubsection{Dritter Schritt: Durchatmen, nachdenken,
planen}

\emph{Was ist wichtig, was kann warten?} -- Weitere Untersuchungen, S-KAMEL und
Maßnahmen.

Nach den Erstmaßnahmen muss man nun planen, welche Maßnahmen,
Untersuchungen und Fragen am dringlichsten sind und zuerst durchgeführt
werden müssen, und was eher in der Abfolge nach hinten ver\-schoben
wird. 

Die Checkliste in Tabelle \ref{tab:checkliste-patientenmanagement} gibt
einen Überblick über die Maßnahmen und das Anamnese\-ge\-spräch, die
Reihenfolge ist bei jedem Patienten anders!

%%%%%%%%%%%%%%%%%%%%%%%%%%%%%%%%%%%%%%%%%%%%%%%%%%%%%%%%%%%

\begin{gWideParEnv}
\begin{minipage}{\linewidth}
\Captionof{table}{Checkliste Patientenmanagement --
Mögliche Maßnahmen, die Reihenfolge ist je nach Patient
unterschiedlich\label{tab:checkliste-patientenmanagement}}

{\sffamily
\begin{tabulary}{\linewidth}[H]{CCL}
\toprule\addlinespace
\Es{Selbstschutz}
%\includegraphics[width=1.005cm,height=0.9cm]{AASdev20090228-img14.gif}
 &
 &
Schutzausrüstung (Handschuhe, Uniform, Helm)
Gefahrenbereich, Unfallstelle absichern, Hunde wegsperren lassen
ggfs. Rückzug, Spezialkräfte anfordern, {\dots}
{\dots}
\\\addlinespace\midrule\addlinespace
\Es{Erkennen}

\includegraphics[width=1cm]{icons/eye.png}
 &
 &
Situation/Milieu {\tiny(zB Alter, soziales Umfeld, SMV,
Alkohol, Pensionistenheim, ...)}

offensichtliche Symptome {\tiny(zB Zyanose,
Bewusstlosigkeit, starke Blutung, ...)}

Schockgefahr {\tiny(Zyanose, kalter Schweiß,
Radialispuls tastbar/nicht tastbar, Puls beschleunigt)}
\\\addlinespace\midrule\addlinespace
\Es{Erste Hilfe}

\includegraphics[width=1cm]{icons/firstAid.png}
 &
 &
Lagerung

Helmabnahme

Blutstillung

psychischer Beistand

Wärmeerhaltung / Kühlung
\\\addlinespace\midrule\addlinespace
{\multirow{6}{2cm}{\centering\includegraphics[width=1cm]{icons/exclamationMark.png}
\Es{Anamnese}}}
 &
{\sffamily\textbf S}
 &
Symptome \& Schmerzen 

{\tiny(Zeit/Verlauf, Stärke, Art/Qualität,
Lokalisation/Ausstrahlung, Provokation)}
\\\addlinespace
 &
{\sffamily\textbf K}
 &
Krankheiten {\tiny(Bestehen Grunderkrankungen, zB
Diabetes?)}
\\\addlinespace
 &
{\sffamily\textbf A}
 &
Allergien/Allergische Reaktion (auch
Allergien auf Medikamente)
\\\addlinespace
 &
 {\sffamily\textbf M}
 &
Medikamente {\tiny(Nimmt Patient Medikamente?
Wogegen sind diese?)}
\\\addlinespace
 &
 {\sffamily\textbf E}
 &
Ereignisse vor Notfalleintritt {\tiny(zB
Anstrengung vor Brustschmerz, Sturz, ...)}
\\\addlinespace
 &
 {\sffamily\textbf L}
 &
Letzte ... (zB letzte Mahlzeit, letzte
Regelblutung, letzter Spitalsaufenthalt)
\\\addlinespace\midrule\addlinespace

{\multirow{8}{2cm}{\centering\includegraphics[width=1cm]{icons/stethoskop.jpg}
\Es{Befunde}}}
 &
{\sffamily\textbf RR}
 &
auskultatorische und palpatorische Messung
\\\addlinespace
 &
{\sffamily\textbf HF}
 &
Herzfrequenz (Puls), Puls rhythmisch oder arrhythmisch
\\\addlinespace
 &
{\sffamily AF}
 &
Atemfrequenz auszählen
\\\addlinespace
 &
{\sffamily Sp\ce{O2}}
 &
Sauerstoffsättigung im Blut, Monitoring
\\\addlinespace
 &
{\sffamily  BZ}
 &
Blutzucker {\tiny(Wann war die letzte
Mahlzeit/Insulinspritze? Diabetes bekannt?)}
\\\addlinespace
 &
{\sffamily Temp.}
 &
Körpertemperatur/Fieber
\\\addlinespace
 &
 {\sffamily\small Neuro\-status}
 &
Bewusstsein (ev. GCS), Orientierung,
Pupillen, Halbseitenzeichen, Herdblick, Meningismus, retrograde Amnesie
\\\addlinespace
 &
{\sffamily\small Trauma\-check}
 &
abnorme Beweglichkeit, DMS an allen
Extremitäten (Nagelbettprobe), Wunden, paradoxe Atmung,
Abwehrspannung des Bauches, ...
\\\addlinespace\midrule\addlinespace
 \Es{Sanitäts\-hilfe}

\includegraphics[width=1cm]{icons/san.png}
&
&
Entscheidung Notarzt-Nachforderung

Fortführung der Ersten Hilfe

Schienung

Sauerstoffgabe

Reanimationsbereitschaft

Verlaufskontrolle

\ldots
\\\addlinespace\midrule\addlinespace
\Es{Notarzt\-assistenz}
\includegraphics[width=1cm]{icons/asb.png}
&
 &
Infusion/periphervenösen Zugang vorbereiten

Intubation vorbereiten

Medikamente vorbereiten

EKG anlegen

\ldots
\\\addlinespace\midrule\addlinespace
\Es{Krankenhaus}

\includegraphics[width=1cm]{icons/hospitalStaff.png}
 &
 &
Welches Spital/Bett buche ich ab? (Arbeitsunfall? Welche Abteilung
fahre ich an? Ist der Patient bereits in einem Spital in Behandlung?)

Brauche ich eine Spezialabteilung? (Schockraum, Stroke Unit, PTCA,
Verbrennungsstation)

Arztübergabe (Anamnese, Befunde, Maßnahmen, Hinweise)
\\\addlinespace\bottomrule
\end{tabulary}}
\end{minipage}
\end{gWideParEnv}

\cite{AMLS3}, \cite{PHTLS6}, \cite{emergency9}, \cite{SemmelABC1}


% \subsection{Der Patient endet nicht mit der Übergabe}
% 
% \begin{quote}
% Aortenaneurysma
% \end{quote}
% 
% \begin{quote}
% Aortenaneurysma 2
% \end{quote}
% 
% \subsection{Fehler und Fallen}
% vdfvrwbwrtbwrtbrwtb
% 
% \begin{quote}
% Ketanest --- {\itshape Eines Tages, es war Winter, die Straßen verschneit und
% glatt, wurden wir zu einer Patientin mit einer offenen Unterschenkelfraktur
% berufen. Die Dame war auf einer nassen Stufe im Hauseingang augerutscht. Sie
% wirkte soweit gefasst, aufgrund des Schocks\footnote{Schock: hier
% psychologischer Schock.} schien sie noch keine großen Schmerzen wahrzunehemn.
% Trotzdem wurde zur prophylaktischen Schmerztherapie ein Notarzt nachbestellt.
% Dieser entschloß sich zu einer Sedierung, u.\,a. mit
% Dormicum{\texttrademark}\footnote{Dormicum{\texttrademark}, Wikstoff Midazolam,
% ist ein Schlafmittel, welches zur Sedierung und Narkose verwendet wird.}. (An
% die anderen Medikemnte kann ich mich nicht mehr genau erinnern, ich denke
% Ketanest{\texttrademark}\footnote{Ketanest{\texttrademark}, Wikstoff S-Ketamin,
% ist ein Narkosemittel.} war auch mit ihm Spiel.) 
% 
% Ich erninnere mich noch dass der Notarzt etwas vom Dormicum{\texttrademark}
% spritzte und sich nach einiger Zeit wunderte weildie erwünschte Wirkung
% ausblieb. Er spritzte noch etwas nach und langsam begann es zu wirken \ldots 
% \emph{"`Na bumm, die hat aber viel gebraucht"'}, merkte er nach der Übergabe
% dann  noch an. Des Rätsels Lösung war aber nicht bei der Patientin zu suchen,
% sondern bei den verwendeten Ampullen: 
% 
% Wir vewendeten Ampullen mit 5\,mg Wirkstoff, die Organisation, der der Notarzt
% angehörte, Ampullen mit 15\,mg (!) Wirkstoff. Daher war es wenig verwunderlich,
% dass "`unser"' Dormicum{\texttrademark} viel weniger Wirkung hatte. Nur gut
% dass der Notarzt "`sein"' Ketanest{\texttrademark} aus der Notarzt-Tasche
% verwendet hatte: Unsere Ketanest{\texttrademark}-Ampullen hatten einen Inhalt von 50\,mg, die des
% Notarztes 25\,mg {\ldots} }
% \end{quote}

\subsection{Sekundäreinsatz}

{\TakeNotesLarge}

\subsection{Ambulanzdienst}

{\TakeNotesLarge}


% \clearpage




\subsection{Übergabegespräch an weiterbehandelndes Fachpersonal}  

\begin{wrapfigure}{r}{4cm}
\includegraphics[width=\linewidth]{./images/shared/aufnahmearzt-ihres-vertrauens-edited2-0500.jpg}
\end{wrapfigure}Die Übergabe des Patienten in den Verantwortungsbereich
eines anderen Sanitäters, eines Notarztes oder des Spitalspersonals hat geordnet
und gewissenhaft zu erfolgen, damit keine wichtigen Informationen verloren
gehen. Diese \E{Schnittstelle} zwischen Sanitäter und Arzt bzw. Pflegepersonal
kann auf den weiteren Behandlungsverlauf einen großen Einfluss nehmen und darf daher in seiner Bedeutung nicht unterschätzt werden! Es empfiehlt
sich -- so Zeit dafür ist -- sich bereits vor Eintreffen des Notarztes oder vor
Ankunft im Krankenhaus zu überlegen, welche Informationen wesentlich sind und
daher auf keinen Fall vergessen werden dürfen.

\vspace{1em}

\noindent Das Übergabegespräch muss mindestens folgende Informationen
beinhalten:
\begin{itemize}
    \item Name des Patienten
    \item Alter
    \item Zusammenfassung der wesentlichen Inhalte der Anamnese
    \item Beschreibung der Symptomatik und deren Verlauf
    \item Bisher erhobene Befunde (Blutdruck, Puls, Sauerstoffsättigung usw.)
    \item Bisher durchgeführte Maßnahmen
    \item Sonstige, wichtige Hinweise (z.B. Infektionsgefahren,
    Bluterkrankheit, Allergien etc.)
\end{itemize}

%=========Ende Abschnitt Emhofer==========================================================

\clearpage
\section{Standardmaßnahmen}

\subsection{Immer durchzuführende Standardmaßnahmen}
\label{topic:standardmassnahmen-immer}

Die folgenden Maßnahmen sind grundsätzlich immer in einer \E{der Situation
angemessenen} Art und Weise durchzuführen. Im begründeten Ausnahmefall kann es
allerdings notwendig oder sinnvoll sein, dass manche Maßnahmen unterbleiben oder
aufgeschoben werden (Selbstschutz, "`Aufklärung"' eines bewusstlosen Patienten,
\ldots{}) oder angepasst werden müssen.

\MassnahmenNG{Standardmaßnahmen}{Immer durchzuführende Standardmaßnahmen}
{\label{m:standardmassnahmen-immer}}{%

\begin{itemize}
  \item Ersteinschätzungsblock.
  \item Erhebung der Anamnese. 
  \item Angemessene Untersuchung. 
  \item Spezielle Maßnahmen gemäß Verdachtsdiagnose(n). 
  \item Situationsgerechte Lagerung. 
  \item Patientenidentifkation.
  \item Dokumentation, Aufklärung.
  \item \E{Wärmeerhalt oder Kühlung.}
  \item Psychischer Beistand.
  \item Verlaufskontrolle, Patientenbeobachtung\,/\,Monitoring.
  
  \item Weiteres Vorgehen, je nach Bedarf und Situation:
  \begin{itemize}
  \item Transport an eine geeignete Einrichtung (Krankenanstalt). 
  \item Notarzt-Nachforderung (bei Bedarf, z.\,B. Schmerztherapie,
  Aufklärung zwecks Belassung auf Patientenwunsch \E{trotz 
  Behandlungsnotwendigkeit}, \ldots).
  \item Be-/Entlassung wenn kein weiterer Behandlungsbedarf besteht und eine
  gesundheitliche Gefährundung des Patienten ausgeschlossen werden kann. 
  \Ca{Rechtliche Hinweise beachten (\prettyref{topic:revers})!}
  \end{itemize}
  \item Ggfs. Übergabe an weiterbehandelndes Personal.
\end{itemize}

{\scriptsize (Die Reihenfolge ist der jeweiligen Situation anzupassen!)}

}

\subsection{Standardmaßnahmen bei vital bedrohten Patienten}
\label{topic:standardmassnahmen-vital}

\MassnahmenNG{Standardmaßnahmen}{\TxMassVit}{\label{m:standardmassnahmen-vital}}{

\begin{enumerate}
  \item \Es{Situationsgerechte Lagerung}
  \item \Es{Sauerstoffgabe}:
  je nach Indikation, Vorsicht bei COPD und Hyper\-ventilation
  \item \Es{Notarzt-Nachforderung}:
  mit kurzer Begründung
  \item \Es{Reanimationsbereitschaft}:
  Platz schaffen, Beatmungsbeutel zu\-sammen\-bauen, ggfs. Defi und
  Absaugung holen und in Griffweite stellen
  \item \Es{Engmaschiges, bestmögliches Monitoring}:
  RR, HF, Pulsoxymetrie wenn vor\-handen, Sitzwache.
\end{enumerate}
}

\clearpage
\section{Fallbeispiele}
\subsection{Primäreinsatz}
\subsubsection{Ein Beispiel: Herr Sepp hat Atemnot}
\Margin{\includegraphics[width=\linewidth]{./images/teaser/rtw-asb-921-alt.jpg}}
\begin{verbatim}
Zeit: dd.04.2003 14:15
Code: 26-B-1 
   Kranke Person (Differentialdiagnose): 
   Unbekannter Zustand (Anrufer 3. Hand) 
BO  : 1210, xxxstraße 13/8/8/45, WHG

\end{verbatim}

\begin{quote}
  \marginlabel{Szeneüberblick}Dia Anfahrt gestaltet sich problemlos. Der
  Berufungsort befindet sich in einem Wohnblock im achten Stockwerk. Das
  Eingangstor wird fixiert um evtl. nachzufordernden Einsatzkräften (Notarzt,
  Polizei, Feuerwehr, \ldots ) den Zugang zu erleichtern. Der Aufzug
  funktioniert und scheint groß genug für die Trage zu sein. 
  
  Sie kommen in eine Wohnung. Die Umgebung erscheint sicher, Sie erkennen auch
  keine Hinweiszeichen auf ein Haustier. Die Wohnung sieht sauber und ordentlich
  aus. Auf der Couch liegt flach ein \marginlabel{Ersteindruck}50-jähriger Mann mit weit aufgerissenen Augen und
  sichtlicher Atemnot. Er wirkt blass und schweißig und fasst sich mit der
  rechten Hand an den Brustbereich.
  
  Sie begrüßen den Patienten, stellen sich vor und
  \marginlabel{Hauptbeschwerde}fragen welche Beschwerden er hat. Der Patient
  \marginlabel{Atemweg, Atmung}atmet schnell und ringt nach Luft, besondere
  Atemgeräusche sind nicht hörbar. Er klagt über Schmerzen in der Brust. Diese
  haben vor einer Viertelstunde angefangen. \marginlabel{Kreislauf}Währenddessen
  haben Sie an der Radialarterie des Patienten den Puls getastet und eine
  Rekap-Probe gemacht. Der Puls ist sehr schwach, schnell und unregelmäßig. Die
  Rekap-Zeit war etwas verlangsamt.
\end{quote}

Welche Informationen haben Sie bis jetzt gewonnen?
\begin{quote}
  Der Patient ist bei Bewusstsein. Die Atemwege scheinen frei, aber er hat
  Atemnot und einen plötzlich aufgetretenen Thoraxschmerz. Er ist blass und
  schweißig, vielleicht steckt ein Schockgeschehen dahinter. Der schwach fühlbare
  Puls und die verlängerte Rekap-Zeit unterstützen diese Vermutung.
\end{quote}

Ist der Patient vital bedroht?
\begin{quote}
  \marginlabel{Vitale Bedrohung: frühzeitig erkannt (noch vor der "`schnellen
  Untersuchung"')}Ja. Zwar wurden noch keine Vitalwerte erhoben, aber bereits
  jetzt gilt der Patient als vital bedroht aufgrund der plötzlich einsetzenden
  Thoraxschmerzen und der Atemnot. Außerdem lässt die Kaltschweißigkeit, der
  schwache Puls und die verlängerte Rekap-Zeit auch nichts gutes vermuten.
\end{quote}

Können Sie schon etwas tun?

\begin{quote}
  \marginlabel{\TxMassVit}Ja.  Bei erhaltenem Bewusstsein, Atemnot und
  Thoraxschmerz wird der Patient mit erhöhtem Oberkörper gelagert. Weiters sprechen
  die Symptome für den Einsatz von Sauerstoff. Es muss noch schnell geklärt werden
  ob eine Hyperventilationstetanie (nicht wahrscheinlich) oder COPD-Erkrankung
  (erfragen) vorliegt.
  
  \TxMassVit:
  
  \begin{itemize}
    \item Der Oberkörper des Patienten wird hochgelagert (Atemnot,
    Thoraxschmerz)
    \item Der Notarzt wird nachgefordert (Berufungsgrund: akuter
    Thoraxschmerz mit Atemnot und schlechtem AZ).
    \item Der Patient wird gefragt ob er an COPD leidet, dies wird verneint.
    Eine psychisch bedingte Hyperventilation erscheint unwahrscheinlich.
    
    Daher wird eine Sauerstoffmaske mit einem Flow von 10\LMin angelegt.
    \item Reanimationsbereitschaft wird hergestellt
    \item Vitalwerte: siehe nächster Schritt.
  \end{itemize}
  
\end{quote}


Was sind die nächsten Schritte?

\begin{quote}
  \marginlabel{Schnelle Untersuchung, Nicht-Traumapatient: Vitalwerte}Wichtig sind
  nun die Vitalwerte und das frühe Beginnen des Monitorings.Dann muss man
  genaueres über das Beschwerdebild in Erfahrung bringen. Der weitere Plan hängt davon ab.
  
  Die Vitalwerte werden gemessen und ein Pulsoxy angehängt (RR
  \RRmmHg{100}{60}, HF 110\,\nicefrac{}{min}, Sp\ce{O2} 94\,\%). Der Patient
  gibt an, dass er stärkste stechende/brennende Schmerzen hinter dem Brustbein hat die ins Kiefer
  ausstrahlen. Außerdem fühle es sich so an als würde jemand auf dem
  Patienten drauf stehen. Die Schmerzen sind nicht atem- oder
  bewegungsabhängig  und haben vor einer Viertelstunde
  {\quotedblbase}wie aus dem Nichts heraus{\textquotedblleft} begonnen.
  Außerdem bekomme er schlecht Luft, sonst hätte er keine Beschwerden.
\end{quote}

\marginlabel{Problem des Patienten}Als Leitsymptome erkennen Sie den
Thoraxschmerz und die Atemnot.

An welche Differentialdiagnosen müssen Sie nun denken?


% \clearpage
% CHECK gWideParEnv Thoraxschmerz: Differentialdiagnosen unseres Patienten
\begin{gWideParEnv}
\begin{tabularx}{\linewidth}[H]{XXXc}
\toprule
\multicolumn{4}{c}{\centering\TabHeading{Thoraxschmerz: Differentialdiagnosen unseres
Patienten} }
\\\midrule
 &
\centering \TabSub{dafür spricht}
 &
\centering \TabSub{dagegen spricht}
 &
\TabSub{Bewertung}
\\\midrule
\emph{Herzinfarkt, Angina pectoris}
 &
Schmerzen hinter dem Brustbein, Ausstrahlung, nicht atemabhängig,
Druck-/Engegefühl, Atemnot
 &
 &
sehr wahrscheinlich
\\\midrule
\emph{Lungenembolie}
 &
Schmerzen hinter dem Brustbein, Atemnot
 &
untypische Ausstrahlung, 
 &
gut möglich
\\\midrule
\emph{Pneumothorax}
 &
Schmerzen im Brustbereich, Atemnot
 &
plötzliche stärkste Schmerzen
 &
eher unwahrscheinlich
\\\midrule
\emph{Trauma}
 &
Schmerzen
 &
kein Hinweis
 &
unwahrscheinlich
\\\midrule
\emph{Erkrankung des Stütz- und Bewegungsapparats}
 &
Schmerzen 
 &
untypische Ausstrahlung, Atemnot, nicht bewegungsabhängig
 &
sehr unwahrscheinlich
\\\bottomrule
\end{tabularx}
\end{gWideParEnv}


\Fazit{Die wahrscheinlichsten Diagnosen sind der
Herzinfarkt und die Lungenembolie.}

Fallen Ihnen bei diesen Erkrankungen wichtige weitere Untersuchungen
ein?

\begin{quote}
  \marginlabel{Durchatmen, nachdenken, planen}Bei beiden Erkrankungen ist es
  wichtig beide Beine zu begutachten. Im Falle einer Lungenembolie sucht man nach
  einer Beinvenenthrombose (einseitig überwärmtes, bläulich-rötliches geschwollenes
  Bein). Bei einer Herzinsuffizienz bzw. Herzinfarkt beurteilt man den Rückstau des
  Blutes, dh. ob die Beine geschwollen sind.
\end{quote}


Außerdem, hat der Patient so etwas schon einmal gehabt? Oder sind einschlägige
Erkrankungen bekannt?  

\begin{quote}
  Sie begutachten beide Beine des Patienten und stellen mittels
  Daumendruckprobe fest, dass beide Beine leicht geschwollen sind. Auf
  Nachfrage gibt der Patient an, dass er sonst keine geschwollenen Beine
  hat.
\end{quote}

Was sind die nächsten Schritte?

\begin{quote}
  Die wichtigsten Untersuchungen wurden gemacht, nun muss man sich um die
  Anamnese kümmern.
  
  Symptome \& Schmerzen: siehe oben.
  
  Krankheiten: Der Patient leidet an nicht-Insulinpflichtigem Diabetes
  mellitus und Bluthochdruck. Sonst sind keine Krankheiten bekannt.
  Spitalsbriefe gibt es nicht.
  
  Sie merken sich, dass Sie nachher den Blutzucker messen wollen!
  
  Allergien: keine bekannt.
  
  Medikamente: Glucophage 850\,mg "`für den Zucker"', Tritace "`für den
  Blutdruck"'. Medikamente wurden heute wie verschrieben eingenommen.
  
  Ereignisse vor Notfalleintritt: Schmerzen begannen plötzlich als der
  Patient Wäsche waschen wollte. In den letzten Wochen immer wieder
  Stiche in der Herzgegend, hat er aber nicht ernst genommen.
  
  Letzte(r):
  
  Spitalsaufenthalt: Vor 10 Jahren in der Rudolfstiftung wegen Rotlauf.
  
  Arztbesuch: Vor einer Woche wegen Neuausstellung eines Rezeptes.
  
  Mahlzeit: Mittagessen, Gulasch
  
  Regel: Der Patient ist männlich!
\end{quote}
% pro:
% 
% Lungenembolie 

Was glauben Sie wie es weiter geht?

\TakeNotesLarge

% ========Subsection eingefügt vom Emhofer 1.5.2009===============================

\section*{Weiterführende Literatur}
\dsliterary{}~~ \cite{SudoweProfHandelnRD1,
AMLS3,
PHTLS6,
emergency9,
SemmelABC1,
KirschnickPflege2,
RueckerBildatlas1,
Ziesenitz2009}

\bibentry{Ziesenitz2009}

\begin{compactdesc}
  \item[\cite{SudoweProfHandelnRD1}] \bibentry{SudoweProfHandelnRD1}
  \item[\cite{AMLS3}] \bibentry{AMLS3}
  \item[\cite{PHTLS6}] \bibentry{PHTLS6}
  \item[\cite{emergency9}] \bibentry{emergency9}
\end{compactdesc}


\texttt{\bibentry{emergency9}}

\textit{\bibentry{emergency9}}

NAEMT: \E{PHTLS. Prehospital Trauma Life Support.} Mosby, 6th edition, September
2006, ISBN 0323033318.

Gernot Rüecker: \E{Bildatlas Notfall- und Rettungsmedizin.}
1. Auflage, 2005, ISBN 3-540-23737-2.

Olaf Kirschnick: \E{Pflegetechniken von A – Z. Schritt für Schritt in
Wort und Bild.} Thieme, 2. Auflage, 2003, ISBN 3-13-127272-4.

Alice Dalton, Daniel Limmer, and Joseph J. Mistovich: \E{Advanced
Medical Life Support: A Practical Approach to Adult Medical Emergencies.} Prentice Hall, 3rd edition, October 2006, ISBN 0131723405.

Aaos: \E{Emergency Care and Transportation of the Sick and Injured.}
Jones and Bartlett Publishers, Inc,      9th edition,    June 2006,
ISBN 0763744050.

Thomas Semmel: \E{ABC – Die Beurteilung von Notfallpatienten.} Elsevier, 1.
Auflage, 2008, ISBN 978-3-437-48560-2.

Hendrik Sudowe: \E{Professionell handeln im Rettungsdienst – Das Trainingsbuch.}
Elsevier, 1. Auflage, 2007, ISBN 978-3-437-48340-0.


